% BEGIN PROFESSOR FEEDBACK PF10: Method - Temporal graph construction
\section{Temporal graph construction}
\label{sec:graph_construction}

The prepared transactions are turned into a temporal graph on which candidate
subgraphs are later counted. This section describes the graph model, how the
graph is actually built from each dataset, and the simple features that are read
straight from it and later serve as a comparison point for the fingerprints.

\subsection{Graph model}

Each address is represented as a node and each transfer as a directed edge that
carries the time at which it happened. The result is a temporal directed
multigraph
\[
  \mathcal{G}=(V,E,\tau,w),
\]
where $V$ is the set of addresses, $E$ the set of transfers, $\tau(e)$ the
timestamp of transfer $e$, and $w(e)$ its transferred value. Two properties are
kept on purpose. Parallel transfers between the same pair of addresses are not
merged, because how many there are and in which order they occur is exactly what
the candidate patterns are built from. Edge direction is kept for the same reason, since an
address that only receives, one that only sends, and one that does both leave
different local role counts behind.

\subsection{Building the graph from the data}

HuggingFace and Real-CATS are stored in very different shapes, so the graph is
assembled differently for each, but always without holding more data in memory
than needed.

For HuggingFace the transaction table is far too large to load as a whole, with
more than 380 million transfers. It is therefore read in small batches. The
first pass measures how active each labelled address is and keeps every transfer
that touches a selected target. These incident transfers define the baseline and
the set containing each target and all of its direct partners. A second full
scan keeps every published transfer whose two endpoints both belong to that set.
The resulting released-neighbourhood closure is used for subgraph counting.

Every triangle involving a selected target consists of two target-to-partner
edges and one edge between those partners. The first two edges place both
partners in the retained node set, and the second scan then keeps their connecting
edge. Triangle counts are therefore complete relative to the published release.
They are not claimed to be complete for Ethereum activity absent from the
release. Baseline features remain restricted to target-incident transfers so
that adding the closure does not give the comparison model extra network scope.

Real-CATS is stored as one transaction history per address, and the same
transfer can appear in the histories of several addresses. These repeated copies
are removed once, using a disk-backed index so that the de-duplication does not
depend on everything fitting in memory. Each history is also split into five
layers, namely ordinary transfers, internal transfers, and ERC-20, ERC-721 and
ERC-1155 token transfers. Every layer becomes its own set of edges, and candidate subgraphs
are counted within each layer. This preserves the distinction between transaction
types, but excludes candidate sequences whose edges cross between layers.

\subsection{The non-motif baseline}

RQ1 asks whether motif fingerprints add anything beyond ordinary address behaviour, so it
needs a description of each address that does not use motifs at all. This
baseline is computed from the same transfers and the same addresses as the motif
fingerprints. Table~\ref{tab:results_baseline_features} lists the resulting 17
features, grouped into how active an address is, how much value it moves, the
structure of its immediate neighbourhood, and how its activity is spread over
time.

\begin{table}[ht]
\centering
\caption{The non-motif baseline features used for RQ1.}
\label{tab:results_baseline_features}
\small
\begin{tabular}{@{}p{3.4cm} p{9.5cm}@{}}
\toprule
\textbf{Feature group} & \textbf{Features} \\
\midrule
Activity &
Number of transfers sent, received, and in total \\[2pt]
Transferred value &
Outgoing and incoming sum, mean, and sample standard deviation, plus the total
value moved \\[2pt]
Neighbourhood structure &
Number of distinct outgoing neighbours, number of distinct incoming neighbours,
and reciprocity, the share of an address's neighbours that it both sends to and
receives from \\[2pt]
Relative timing &
Active span and the mean, median, and sample standard deviation of the gaps
between successive transfers that touch the address \\
\bottomrule
\end{tabular}
\end{table}

A few choices keep this baseline honest. Absolute first and last transaction
times are deliberately left out, because they mostly reveal when the data was
collected or how old a campaign is rather than how an address behaves, and a
model could otherwise separate the classes by collection window instead of by
behaviour. Measures that describe an address's position in the wider network are also left
out, because computing them reliably would need more of the surrounding graph
than is kept here. Where a value cannot be computed, for instance a measure of
variation for an address with too few transactions, it is set to zero.
% END PROFESSOR FEEDBACK PF10

